% ---------------------------------------------------------------------------
% Appendix: Migration from earlier versions
% ---------------------------------------------------------------------------
\chapter{Migration from Earlier Versions}\label{app:migration}
% Long monospaced identifiers leave little room for line breaking; allow some
% extra inter-word stretch in this chapter (the group keeps it local).
\begingroup\setlength{\emergencystretch}{3em}% chapter-local

This appendix explains what code written for earlier versions of the package
has to change. Section~\ref{sec:migration-11} covers the step from version 1.0
to 1.1, which adds features and changes little existing behaviour.
Section~\ref{sec:migration-pre10} covers code written for the versions before
1.0, when the package consisted of two separate projects. The complete list of
changes is in \code{CHANGELOG.md}.

\section{From 1.0 to 1.1}\label{sec:migration-11}

Version 1.1 adds the environments \envid{PowerGrid-v0} and \envid{Platoon-v0},
native vector environments, the catalogue and the \code{env-lib} command line,
wrappers, baseline controllers, evaluation tools and graph utilities
(Chapter~\ref{chap:workflow}). No public name was removed, and class names,
module paths, constructor arguments and their defaults are unchanged. The
following changes can affect code written for version 1.0.

\paragraph{Gymnasium 1.0 or newer.} The package requires
\code{gymnasium>=1.0} (version 1.0 accepted 0.29). The vector API of
Gymnasium changed in 1.0: a finished copy of a vector environment is by
default reset on the \emph{next} call to \code{step()}, whereas Gymnasium 0.29
reset it within the same step, and in same-step mode the last observation is
returned in \code{infos["final\_obs"]} instead of Gymnasium 0.29's
\code{infos["final\_observation"]}. Code that ran vector
environments with Gymnasium 0.29 has to follow these conventions
(Section~\ref{sec:workflow-vector}); single environments are not affected.

\paragraph{Unknown reset options warn instead of raising.} In version 1.0 the
Consensus, Formation and Kuramoto environments raised \code{ValueError} for a
key of \code{reset(options=...)} that they did not know. They now ignore such
keys with a \code{UserWarning}, as do the new PowerGrid and Platoon
environments, so that generic tools can pass options of their own; the values
of known keys are still validated and raise \code{ValueError} when invalid
(Section~\ref{sec:envlib-reset-options}). Code or tests that relied on the
error to catch a misspelt key can turn the warning into an error:

\begin{pycode}
import warnings
import numpy as np
import env_lib

env = env_lib.make("Consensus-v0")
with warnings.catch_warnings():
    warnings.filterwarnings("error", message=r".*unknown reset option")
    try:
        env.reset(seed=0, options={"position": np.zeros((8, 2))})   # misspelt
    except UserWarning as err:
        print(err)
\end{pycode}

\paragraph{\texorpdfstring{\code{max\_episode\_steps}}{max\_episode\_steps} sets the environment's own limit.}
In version 1.0, \code{env\_lib.make(id, max\_episode\_steps=N)} passed the
argument on to \code{gymnasium.make}, which added a \code{TimeLimit} wrapper on
top of the environment's own limit: the shorter of the two applied, so
\code{env\_lib.make("Consensus-v0", max\_episode\_steps=500)} still ended
episodes after 200 steps, and for AJLATT the value never reached the
configuration. \code{env\_lib.make} and \code{env\_lib.make\_vec} now set the
environment's limit argument (\code{max\_steps}, \code{max\_iter},
\code{max\_cycles} or \code{max\_episode\_steps}; see
\code{env\_lib.get\_spec(id).limit\_kwarg}) instead, so the value given is the
episode length and no \code{TimeLimit} wrapper is added. Consequently
\code{env.spec.max\_episode\_steps} is \code{None}; read the limit from the
environment, for example \code{env.unwrapped.max\_steps}. Passing both
\code{max\_episode\_steps} and a different value of the limit argument raises
\code{ValueError}. \code{gymnasium.make} called directly behaves as before.

\paragraph{Native vector environments.} The ids \envid{Consensus-v0} and
\envid{Formation-v0} and the NumPy Kuramoto ids now register a vector entry
point, as do the new PowerGrid and Platoon ids. For these ids,
\code{gymnasium.make\_vec(id, n)}, which returned a \code{SyncVectorEnv} in
version 1.0, and \code{env\_lib.make\_vec(id, n)} therefore return the native
\code{ConsensusVectorEnv} or \code{KuramotoOscillatorVectorEnv}. They follow the
same dynamics and the Gymnasium conventions of \code{SyncVectorEnv} (next-step
autoreset, where the reset step ignores the copy's action and returns reward
0, both flags \code{False} and the info of the reset; batched infos with
\code{"\_key"} masks), but differ in the following points:
\begin{itemize}
  \item All copies draw from one generator: copy $i$ no longer reproduces a
    single environment reset with seed $s + i$
    (Section~\ref{sec:workflow-vector}).
  \item There is no \code{envs.envs} list of single environments; the batched
    state is available through properties such as \code{envs.positions}
    (Appendix~\ref{app:api}).
  \item \code{render()} returns the frame of copy 0 instead of a tuple with
    one frame per copy.
  \item The autoreset mode is a keyword argument of \code{make\_vec}
    (\code{autoreset\_mode="same\_step"}, for every vectorisation mode), not an
    entry of \code{vector\_kwargs}.
\end{itemize}
Pass \code{vectorization\_mode="sync"} to keep the \code{SyncVectorEnv} of
version 1.0.

\paragraph{AJLATT example controller.} The function \code{encircle\_policy(env,
radius, gain)} of \code{examples/ajlatt\_demo.py} was removed from the example.
Its controller is part of the package as \code{env\_lib.baselines.ajlatt\_encircle}
(a pure function of the observation) and is returned, with its parameters bound
to the environment, by \code{env\_lib.baseline\_policy(env)}:

\begin{pycode}
import env_lib
from env_lib import AJLATTEnv

env = AJLATTEnv(num_robots=4)
# Before (examples/ajlatt_demo.py of version 1.0): action = encircle_policy(env)
policy = env_lib.baseline_policy(env, radius=1.8, gain=1.2)   # the old defaults
obs, info = env.reset(seed=0)
obs, reward, terminated, truncated, info = env.step(policy(obs))
\end{pycode}

The old function read the robots' pose and target estimates from the
environment's internal filters; the packaged controller reads the same
quantities from the observation, in the robot's frame. Its results are
therefore close to, but not identical with, those of version 1.0: the two
default episodes of \code{ajlatt\_demo.py} return $-7333.0$ and $-7113.1$
instead of $-7334.3$ and $-7170.4$.

\paragraph{The catalogue module is callable.} \pkg{env\_lib.catalog} is both a
module and a function: \code{env\_lib.catalog(...)} returns the catalogue
whether or not the submodule has been imported, and its other names remain
module attributes (\code{env\_lib.catalog.describe},
\code{env\_lib.catalog.EnvInfo}, \code{env\_lib.catalog.Catalog}).
\code{from env\_lib.catalog import catalog} imports the function itself.

\section{From versions before 1.0}\label{sec:migration-pre10}

Before version 1.0 the code lived in two separate projects,
\code{envlib\_project/env\_lib} and \code{toolkit\_project/toolkit}, each with
its own \code{setup.py}, and several environments had their own installation
files. Version 1.0 turns them into one distribution with a stable API, fixes a
number of bugs and changes some behaviour. This section explains, module by
module, what code written for these earlier versions has to change.

Most user code keeps working: class names, module paths, constructor
arguments and their defaults are unchanged, and deprecated arguments are
translated with a \code{DeprecationWarning}. Python shows these warnings only
for code in \code{\_\_main\_\_} by default; run your scripts once with
\code{python -W default::DeprecationWarning} to list everything that needs
attention.

\subsection{Checklist}\label{sec:migration-checklist}

\begin{enumerate}
  \item Uninstall the old distributions and install \pkg{my-tool-box}
    (Section~\ref{sec:migration-install}).
  \item Unpack \code{(observation, info)} from every \code{reset()} call and use
    the five-tuple of \code{step()}, including for AJLATT.
  \item Create AJLATT through \code{env\_lib.AJLATTEnv} with the new parameter
    names and adapt to the joint action and observation spaces
    (Section~\ref{sec:migration-ajlatt}).
  \item For Pistonball in discrete mode, pass one value in $\{0, 1, 2\}$ per
    piston (Section~\ref{sec:migration-pistonball}).
  \item Treat \code{terminated} and \code{truncated} separately, and set episode
    limits through the constructor arguments
    (Section~\ref{sec:migration-common}).
  \item Retrain \pkg{neural\_toolkit} models; old checkpoints cannot be loaded
    (Section~\ref{sec:migration-neural}).
  \item If your figures relied on side effects of \pkg{plotkit} on the global
    matplotlib configuration, call \code{set\_research\_style()} explicitly
    (Section~\ref{sec:migration-plotkit}).
\end{enumerate}

\subsection{Installation}\label{sec:migration-install}

\textbf{Before}, each project (and optionally each environment) was installed
from its own folder:

\begin{shellcode}
cd envlib_project && pip install -e .
cd ../toolkit_project && pip install -e .
cd ../envlib_project/env_lib/pistonball_env && pip install -e .   # and similar
\end{shellcode}

\textbf{After}, one installation from the repository root provides both
import packages; optional features are selected with extras
(Chapter~\ref{chap:install}):

\begin{shellcode}
pip uninstall -y env_lib toolkit
pip uninstall -y ttenv kos_env linemsg_env pistonball_env wireless_comm_env
pip install -e ".[all]"            # or ".[all,dev]" for development
\end{shellcode}

The import names \pkg{env\_lib} and \pkg{toolkit} are unchanged. TensorFlow,
seaborn, pandas and scikit-learn are no longer required. Importing
\pkg{env\_lib} or \pkg{toolkit} no longer imports PyTorch, pygame, pymunk or
Tk; the subpackages and environment classes are loaded on first use, so
\code{import toolkit; toolkit.plotkit.plot\_line(...)} and
\code{from env\_lib import PistonballEnv} continue to work.

\subsection{Changes common to all environments}\label{sec:migration-common}

\paragraph{Reset and step.}
All environments follow the Gymnasium API. \code{reset(*, seed=None,
options=None)} returns \code{(observation, info)}, and
\code{step(action)} returns \code{(observation, reward, terminated,
truncated, info)}. The \code{seed} argument must be passed by keyword.

% norun
\begin{pycode}
# Before (AJLATT and old-style loops)
# obs = env.reset()
# obs, reward, done, info = env.step(action)

# After
obs, info = env.reset(seed=0)
obs, reward, terminated, truncated, info = env.step(env.action_space.sample())
done = bool(np.any(terminated)) or truncated
\end{pycode}

\paragraph{Seeding.}
\code{env.seed(s)} is deprecated where it existed (LineMsg, WirelessComm,
Pistonball) and emits a \code{DeprecationWarning}; use
\code{env.reset(seed=s)}. No environment seeds the global NumPy or PyTorch
generators any more (the Kuramoto environments used to call
\code{np.random.seed(42)} and \code{torch.manual\_seed(42)}); code that relied
on these side effects for its own random numbers must seed its own generator.

\paragraph{Episode limits and termination.}
The registered ids no longer add a \code{TimeLimit} wrapper. Previously, for
example, \envid{KuramotoOscillator-v0} was truncated after 1000 steps by the
wrapper even if a larger \code{max\_steps} was passed to \code{gym.make}; now
the constructor argument (\code{max\_steps}, \code{max\_iter},
\code{max\_cycles} or \code{max\_episode\_steps}) is the only limit, and
reaching it sets \code{truncated=True}. Task success (for example
synchronisation of the oscillators) sets \code{terminated=True}; previously
some environments reported both cases as \code{terminated}.

\paragraph{Registration.}
Environments are registered when \pkg{env\_lib} is imported; the
\code{register()} functions of \pkg{env\_lib.kos\_env} and
\pkg{env\_lib.wireless\_comm\_env} are deprecated no-ops that call the central
registration. \code{env\_lib.make(id)} and \code{env\_lib.list\_envs()} are
new. The id \envid{AJLATT-v1}, an unusable alias of \envid{AJLATT-v0}, was
removed.

\paragraph{Observations, rewards and rendering.}
Observations are \code{float32} and lie in \code{observation\_space};
per-agent rewards are available in \code{info["agent\_rewards"]}. Rendering
happens only on request: in \code{"human"} mode the environment renders
itself in \code{reset()} and \code{step()}, so explicit \code{env.render()}
calls in \code{"human"} loops are no longer needed, and \code{"rgb\_array"}
works without a display. \code{step()} before the first \code{reset()} now
raises an error instead of resetting implicitly.

\subsection{Kuramoto oscillators}\label{sec:migration-kuramoto}

\begin{itemize}
  \item \emph{PyTorch backend.} A sign error made the coupling repulsive
    ($\sin(\theta_i - \theta_j)$ instead of $\sin(\theta_j - \theta_i)$).
    Results obtained with \code{KuramotoOscillatorEnvTorch} before 1.0 are
    therefore not comparable with new ones. Both backends now produce the same
    trajectories.
  \item \emph{Normalisation.} Division of the coupling sum by $N$ is opt-in
    for both backends (\code{normalize\_coupling=True}). Code that relied on
    the normalisation of the old PyTorch backend must set this flag.
  \item \emph{Coupling matrix.} The NumPy backend used to skip entries
    $K_{ij} \le 0$; the full matrix is used now. Default settings are not
    affected.
  \item \emph{Constant coupling.} \code{coupling\_mode="constant"} without a
    matrix now uses \code{coupling\_strength * adjacency}, which makes the
    registered \envid{*-Constant-v0} ids usable (they raised before).
  \item \emph{Episode end.} Synchronisation ($r >$ \code{sync\_threshold})
    sets \code{terminated}; \code{sync\_threshold} and \code{sync\_bonus} are
    constructor arguments. The PyTorch backend accepts \code{device="auto"},
    and its phase coherence and bonus rules now match the NumPy backend.
\end{itemize}

\subsection{LineMsg and WirelessComm}\label{sec:migration-networked}

Seeded trajectories are bit-identical to the previous implementations, so
existing results remain valid. The text output of the old renderers is
available as \code{render\_mode="ansi"}; \code{"human"} and \code{"rgb\_array"}
show the new dashboards. LineMsg accepts actions either as an integer in
\code{Discrete(2**N)} or as an array of $N$ bits, and
\code{action\_space\_type="multibinary"} supports more than 62 agents.
WirelessComm validates its arguments and actions and reports per-agent
outcomes (idle, success, collision, lost) in \code{info["outcomes"]}.

\subsection{Pistonball}\label{sec:migration-pistonball}

\paragraph{Discrete actions.}
The discrete action space is \code{MultiDiscrete([3] * n\_pistons)}: one
value per piston, 0 moves the piston down, 1 keeps it still and 2 moves it up.
Earlier versions declared the same array format as
\code{Box(0, 2, (n\_pistons,), int32)}, and some earlier documentation used a
single base-3 integer in \code{Discrete(3**n\_pistons)}. Code that sends
arrays keeps working; code that encodes a single integer must decode it:

\begin{pycode}
import numpy as np
import env_lib

env = env_lib.PistonballEnv(n_pistons=5, continuous=False)
obs, info = env.reset(seed=0)

# Before: one integer for all pistons (base-3 encoding)
code = 13
# After: one value per piston; decode an old integer action like this
action = np.array([(code // 3**i) % 3 for i in range(env.n_pistons)])  # [1 1 1 0 0]
obs, reward, terminated, truncated, info = env.step(action)
\end{pycode}

\paragraph{Other changes.}
Actions are validated (shape, range, NaN). The screen is only drawn when
rendering is requested; the methods \code{draw()}, \code{draw\_background()},
\code{draw\_pistons()} and \code{enable\_render()} are deprecated no-ops, and
the PNG sprites were replaced by a vector renderer. The environment works with
pymunk 6 and 7.

\subsection{AJLATT}\label{sec:migration-ajlatt}

AJLATT was rewritten as a Gymnasium environment configured by a dataclass.
The old factory called \code{argparse.parse\_args()} internally, which failed
in any script with its own command-line arguments.

\paragraph{Creating the environment.} The legacy factory still works; the
class and the configuration dataclass are the preferred interface:

\begin{pycode}
import env_lib
from env_lib import AJLATTConfig, AJLATTEnv

# Before (still works, with DeprecationWarnings for the legacy names)
env = env_lib.ajlatt_env(map_name="obstacles04", num_Robot=4, T_steps=100)

# After
env = AJLATTEnv(map_name="obstacles04", num_robots=4, max_episode_steps=100)

# Or with an explicit configuration object
config = AJLATTConfig(map_name="obstacles05", num_robots=3)
env = AJLATTEnv(config, render_mode="rgb_array")
\end{pycode}

Table~\ref{tab:trouble-ajlatt} lists the renamed parameters. Rendering is off
by default; the old \code{render=True} flag becomes
\code{render\_mode="human"}, and \code{render\_mode="rgb\_array"} records
frames without a display.

\paragraph{Spaces, rewards and termination.}
The spaces are \emph{joint}: \code{action\_space} has shape
\code{(num\_robots, 2)} (linear and angular velocity of each robot) and
\code{observation\_space} shape \code{(num\_robots, obs\_dim)}, both
\code{float32}. The per-robot spaces used by the old version are available as
\code{single\_action\_space} and \code{single\_observation\_space}. The reward
and \code{terminated} are per-robot arrays; wrap the environment in
\code{TeamRewardWrapper} for a scalar team reward and termination flag.

\begin{pycode}
import numpy as np
from env_lib import AJLATTEnv
from env_lib.ajlatt_env import TeamRewardWrapper

env = AJLATTEnv(num_robots=4)
print(env.single_action_space.shape, env.action_space.shape)   # (2,) (4, 2)

obs, info = env.reset(seed=0)                    # before: obs = env.reset()
actions = np.stack([env.single_action_space.sample() for _ in range(4)])
obs, reward, terminated, truncated, info = env.step(actions)
print(reward.shape, terminated.shape)            # (4,) (4,)

team_env = TeamRewardWrapper(AJLATTEnv(num_robots=4))
obs, info = team_env.reset(seed=0)
obs, team_reward, done, truncated, info = team_env.step(actions)   # float, bool
\end{pycode}

\paragraph{Behaviour changes.} The following changes affect results:
\begin{itemize}
  \item The target-velocity entries of the observation contain the commanded
    target velocity instead of the configured constant.
  \item Covariance intersection uses a Newton solver by default; its objective
    is never worse than that of the previous SLSQP solver. Pass
    \code{ci\_solver="slsqp"} to reproduce results of earlier versions (to
    about $10^{-6}$).
  \item Target policies are per environment instance and selected with
    \code{target\_policy}; they used to be shared class attributes that were
    not reset between episodes. Unknown map names raise an error instead of
    prompting with \code{input()}.
  \item Covariance traces were appended to lists that were never cleared (a
    memory leak). Per-episode traces are returned by
    \code{env.episode\_statistics()}.
  \item Dynamic maps are regenerated at every \code{reset()} and no longer
    need scikit-image.
\end{itemize}

\paragraph{Training scripts.}
The PPO, LLM and debugging options that used to live in the AJLATT parameter
module were removed from the environment package. To expose the environment
options on a script's own parser, use \code{AJLATTConfig.add\_arguments}:

\begin{pycode}
import argparse
from env_lib import AJLATTConfig, AJLATTEnv

parser = argparse.ArgumentParser()
parser.add_argument("--lr", type=float, default=3e-4)       # your own options
AJLATTConfig.add_arguments(parser, prefix="env.")           # --env.num_robots, ...
args = parser.parse_args(["--env.num_robots", "3"])
config = AJLATTConfig.from_namespace(args, prefix="env.")
env = AJLATTEnv(config)
\end{pycode}

\paragraph{Maps.}
The map tools moved to \pkg{env\_lib.ajlatt\_env.maps} and are available as
the console scripts \code{ajlatt-build-map} and \code{ajlatt-plot-map}
instead of the scripts \code{build\_map.py} and \code{plot\_map.py}. The old
builder wrote transposed grids; rebuild custom maps from their specification
with the new builder. Maps whose grid is one row shorter than the header
declares are padded with a wall row and a warning (three bundled maps had this
problem and were fixed).

\subsection{plotkit}\label{sec:migration-plotkit}

The function names and positional arguments are unchanged; new options are
keyword-only arguments appended to the signatures.

\begin{itemize}
  \item \emph{No global side effects.} The plot functions used to call
    \code{plt.style.use("default")} and change \code{rcParams}, which also
    affected later, unrelated figures. They now apply their style locally.
    Code that relied on the side effect should call \code{set\_research\_style()}
    once (and \code{reset\_style()} to undo it). \code{set\_research\_style} no
    longer sets \code{figure.dpi=300}, which made on-screen figures very large;
    it still sets \code{savefig.dpi=300}.
  \item \emph{Fixed input handling.} \code{plot\_line(x, [y1, y2])} plots both
    series (it plotted only the first), and a flat list of numbers is one
    series (it was treated as several). Several bar series are grouped side by
    side instead of overlapping.
  \item \emph{New features.} \code{band} and \code{smoothing} in
    \code{plot\_shadow\_curve}, NaN-robust statistics,
    \code{plot\_learning\_curves}, \code{plot\_histogram}, \code{save\_figure},
    style presets, the Okabe-Ito palette and the gallery command. The heatmap
    functions no longer need seaborn.
\end{itemize}

\begin{pycode}
from toolkit.plotkit import plot_line, reset_style, set_research_style

# Before: calling any plotkit function changed the global style as a side effect.
# After: opt in explicitly when plain matplotlib code should use the research style.
set_research_style()
ax = plot_line([0, 1, 2], [[0, 1, 4], [0, 2, 3]], labels=["a", "b"])   # two series
reset_style()
\end{pycode}

\subsection{neural\_toolkit}\label{sec:migration-neural}

\paragraph{Retrain models.}
Earlier MLP stacks applied no activation between hidden layers, and several
other defects changed the computed functions (Section~\ref{sec:neural-changes}).
The layer structure of every network changed, so checkpoints written before
1.0 cannot be loaded; retrain and save new checkpoints.

\paragraph{API differences.} Besides the fixes, the following details changed:
\begin{itemize}
  \item Unknown activation names raise \code{ValueError} instead of silently
    falling back to a default, and the \code{device} argument is honoured
    (modules are moved to it at construction).
  \item Q-networks honour the \code{action} argument: \code{q(s, a)} returns
    $Q(s, a)$ with shape \code{(B,)} for integer actions.
  \item \code{NetworkUtils.load\_checkpoint} loads with
    \code{weights\_only=True} and maps tensors to the model's device by
    default. Checkpoints that contain arbitrary Python objects in the extra
    entries need \code{weights\_only=False} (only for trusted files).
    \code{save\_checkpoint} accepts \code{optimizer=None} and extra keyword
    entries, and returns the path.
  \item The method \code{create\_positional\_encoding} of \code{NetworkUtils}
    returns the frozen table as an \code{nn.Parameter} of shape
    \code{(1, max\_len, d\_model)};
    use \code{PositionalEncoding} for a module that adds it to its input.
  \item The tabular tools draw random numbers from a per-table generator
    (\code{rng} argument) instead of the global NumPy state.
\end{itemize}

\paragraph{Names from earlier documentation.}
The documentation of earlier versions described some names that never existed
in the code. Table~\ref{tab:migration-neural-names} lists the actual
equivalents.

\begin{table}[htbp]
  \centering
  \small
  \caption{Names used in pre-1.0 documentation and their actual equivalents.}
  \label{tab:migration-neural-names}
  \begin{tabularx}{\textwidth}{@{}>{\raggedright\arraybackslash}p{5.6cm}>{\raggedright\arraybackslash}X@{}}
    \toprule
    Documented before 1.0 & Use \\
    \midrule
    \multicolumn{2}{@{}l}{\emph{Static methods of} \code{NetworkUtils}} \\
    \code{count\_trainable\_parameters} & \code{count\_parameters} (counts trainable parameters) \\
    \code{init\_weights} & \code{initialize\_weights} \\
    \code{clip\_gradients} & \code{clip\_grad\_norm} \\
    \code{get\_device}, \code{to\_device} & plain PyTorch: \code{torch.device(...)}, \code{module.to(device)} \\
    \midrule
    \multicolumn{2}{@{}l}{\emph{Encoder arguments}} \\
    \code{output\_dim} of \code{CNNEncoder} and \code{RNNEncoder} & \code{latent\_dim} \\
    \bottomrule
  \end{tabularx}
\end{table}

\subsection{parakit}\label{sec:migration-parakit}

\code{ParameterTuner} (and its alias \code{ParameterAdjuster}) keeps its
constructor and \code{tune\_parameters} class method, but its behaviour was
corrected:

\begin{itemize}
  \item \code{tune()} installs the accepted values as the parser's defaults, as
    it was documented to do; previously the returned parser was unchanged.
  \item Boolean strings (the text \code{"False"} used to be stored as a
    string, which is true in a boolean context), lists (\code{nargs}),
    \code{None} values and \code{store\_true} options are parsed correctly.
  \item Validation callbacks receive the \emph{converted} value. Callbacks
    written for the old convention, which received text and returned
    \code{(is\_valid, converted\_value, message)}, keep working as long as they
    also accept an already converted value (for example \code{float(value)} of
    a float); their converted value is ignored.
  \item \code{import toolkit.parakit} works without Tk; the window imports Tk
    only when it opens. The library no longer calls
    \code{logging.basicConfig}, so configure logging in your script if you want
    to see its log messages.
  \item New headless functions (\code{get\_parameters},
    \code{save\_parameters}, \code{load\_parameters}, \code{apply\_parameters},
    \code{parse\_value}) make the GUI optional
    (Section~\ref{sec:parakit}).
\end{itemize}

\begin{pycode}
import argparse
from toolkit.parakit import apply_parameters, load_parameters, save_parameters

parser = argparse.ArgumentParser()
parser.add_argument("--learning_rate", type=float, default=1e-3)

# Before: validators received text and returned a 3-tuple
def validate_lr_old(value):
    lr = float(value)
    return (0 < lr <= 1), lr, "learning rate must be in (0, 1]"

# After: validators receive the converted value; a bool or (ok, message) suffices
def validate_lr(value):
    return 0 < value <= 1, "learning rate must be in (0, 1]"

path = save_parameters({"learning_rate": 0.01}, "runs/params")
apply_parameters(parser, load_parameters(path),
                 validation_callbacks={"learning_rate": validate_lr_old})
apply_parameters(parser, {"learning_rate": "0.02"},
                 validation_callbacks={"learning_rate": validate_lr})
\end{pycode}

\par\endgroup
